
This paper reported existence-tier evidence for depth-resolved, training-free hallucination detection. Across three 7--8B model families and two QA datasets --- 5,451 scored generations covering all six model $\times$ dataset cells --- at least one screened intermediate layer, always in the L28--L31 band just before output calibration, is class-separable using raw logit-lens uncertainty statistics read from one greedy generation plus one teacher-forced re-forward (selection-split corrected AUROC 0.6092--0.7011 against a 0.55 gate), and the best such layer exceeds the same sweep's final-layer entropy readout in all six cells as point estimates. No probe was trained and no sample repeated. Adjacent-layer KL divergence, included as the third signal, is --- to our knowledge --- used here for the first time as a detection-time score; it is the best signal in two of the six cells and provides the top three intermediate scores in the binding LLaMA-2/TriviaQA cell.

The second finding concerned measurement. The documented-failure baseline this study set out to compare against did not survive its own provenance audit: six ordinary protocol differences moved the same cell's final-layer AUROC from 0.5186 to 0.5928, and a designed halt gate diagnosed the unsatisfiable anchor in 47 seconds. The protocol-internal anchor that replaced it --- identity-verified cache reuse, within-sweep baselines, descriptive-only cross-run comparisons --- is a methodological contribution that arose directly from that failure.

All evidence reported here is selection-split point estimates from a single seed, with no confidence intervals; the four dependent verification tiers (test-split confirmation, cross-dataset transfer, KL--entropy fusion, no-regression guard) were never executed and their predictions remain inconclusive. The staged design makes the next tier inexpensive: the highest-value follow-ups --- freezing the six selected tuples and evaluating the locked test splits with a paired bootstrap, the transfer matrix, the fusion test, and confidence intervals on the existence grid --- cost zero GPU from the released caches. Two follow-ups need modest compute: base-vs-instruct pairs of one family, to test the reversed tuning-polarity hypothesis suggested by the +0.130 instruct-model margin, and a factorial attribution of the six protocol factors behind the 0.5186 $\rightarrow$ 0.5928 shift. Within the tested scope, the results support a simple summary: the output layer is where the model speaks, but it is not the most informative place to read its uncertainty.
